\begin{theorem}[Soundness and completeness of linear temporal logic]
For every temporal formula \(A\),
\[
  \derives_{\mathsf{LTL}}A
  \quad\Longleftrightarrow\quad
  \mathcal M,n\entails A
  \text{ for every temporal structure \(\mathcal M\) and every \(n\in\N\)}.
\]
This is weak completeness for validity, not a strong-completeness assertion
for arbitrary premise sets.  It concerns discrete infinite linear time;
finite traces, dense time, and branching time require different semantics and
calculi.
\end{theorem}

\paragraph{Proof sketch.}
Soundness of the first two temporal schemas and of
\textsc{XNec} follows directly from the unique successor.  The unfolding
schema follows by separating the cases where the witness for until is the
current position or a later position.  For \textsc{UInd}, assume its premise is
valid and reason backward from the least witness for \(B\) to show \(C\) at
each preceding position at which \(A\) holds.  Completeness closes a formula
under subformulas, negation, and the temporal obligations generated by until;
it constructs maximally consistent closure sets linked by the \(\mathsf X\)
requirements, proves that every pending until obligation is eventually
realized, and unravels the resulting path into a temporal structure.  The
truth
lemma then yields a falsifying temporal structure for every underivable
formula.  The consistency, eventuality-realization, and truth-lemma case
calculations are omitted---this is why the theorem is presented with a
proof sketch rather than a complete proof---and are controlled by
\cite{GabbayPnueliShelahStavi1980}.

